\begin{textsample}{POS Dim 1 – vem\_ed\_16 – Score 11.00 – vem\_ed\_16/t069.txt}  \label{ex:f1_pos_090}
QUEM \textbf{É} QUEM Márcio Venício Barbosa \textbf{é} presidente da Associação Brasileira de Educação Internacional (FAUBAI), \textbf{que} atua na ampliação do \textbf{processo} de internacionalização das Instituições de Ensino Superior (IES) brasileiras, por meio de projetos como o BRaVE – Brazilian Virtual Exchange, \textbf{programa} \textbf{que} incentiva a implantação de Intercâmbios Virtuais.
Professor titular de Literatura Francesa e Secretário de Relações Internacionais da Universidade Federal do Rio Grande do Norte (UFRN), Barbosa concedeu, em entrevista por e-mail, suas \textbf{reflexões} sobre Internacionalização em Casa.
Comente a \textbf{importância} da Internacionalização em Casa para as instituições de ensino \textbf{superior} no Brasil.
Vivemos em um país de proporções continentais e, não há como negar, nós mesmos desconhecemos algumas regiões de nosso território ou, se conhecemos, \textbf{é} sem muita profundidade.
Essa peculiaridade geográfica, associada à hegemonia da língua portuguesa na quase totalidade do território nacional, nos \textbf{torna} um tanto fechados, não ao contato com estrangeiros, pois os temos em grande \textbf{número}, mas fechados ao contato mais estreito com suas culturas, suas línguas e suas visões de mundo.
Recebemos muito bem os estrangeiros, desde \textbf{que} estejam aculturados, \textbf{que} se expressem, mesmo mal, em português e não \textbf{representem} um motivo de preocupação, como \textbf{é} o caso, infelizmente, dos refugiados...
Num cenário como esse, não basta \textbf{que} tenhamos uns poucos alunos brasileiros estudando no exterior para \textbf{que} sejamos internacionalizados.
Precisamos também ter uma maior presença internacional em nossas IES, como \textbf{forma} de trazer o mundo até nós.
Dificilmente poderemos lançar e custear \textbf{programas} \textbf{que} atendam milhares de estudantes em experiências de mobilidade presencial.
Por isso \textbf{é} importante destinar recursos para atividades menos onerosas e capazes de atingir um \textbf{número} maior de estudantes.
Projetos de cooperação \textbf{que} tragam estudantes e pesquisadores internacionais a nossas instituições \textbf{são} um excelente vetor para a Internacionalização em Casa, pois isso \textbf{permite} uma visibilidade maior da instituição e pode preparar melhor \textbf{toda} a comunidade acadêmica para o convívio com outras culturas.
Precisamos de mais investimento tanto na infraestrutura de nossos espaços de ensino-aprendizagem, \textbf{quanto} na formulação de nossas \textbf{propostas} de ensino, como a internacionalização dos currículos, passando, também, pela capacitação de nossos docentes para o ensino dos \textbf{conteúdos} de suas respectivas áreas em línguas estrangeiras.
Essas práticas poderão assegurar um \textbf{processo} institucionalizado de internacionalização, abrindo as portas do mundo para \textbf{toda} a comunidade acadêmica.

% matched lemmas: conteúdo, forma, importância, número, permitir, processo, programa, proposta, quanto, que, reflexão, representar, ser, superior, todo, tornar
\end{textsample}
